% SHARD-ID: AQM-84-QSA-INFINITE-RING-BOUNDARIES
% SHARD-TITLE: Infinite Ring Boundary Cases
% SHARD-SUMMARY: Reviews \(\Spec\mathbb F_q[t]\), \(\Spec\mathbb F_q[x,y]\), and \(\Spec\mathbb Z\) as infinite QSA boundary examples.
% SHARD-SUMMARY: Separates finite supports, algebraic quasi-local algebras, reference-vector tensor representations, and formal infinite-volume profiles.
% SHARD-SUMMARY: Records missing representation and completion conventions before the Step 21 agreement and inequivalence synthesis.
% SHARD-KEYWORDS: quantum-system association, infinite rings, quasi-local algebra, incomplete tensor product, generic point, completion gap

\section{Infinite Ring Boundary Cases}
\label{sec:qsa-infinite-ring-boundaries}

\subsection{Status and Local Anchors}

\textbf{Status: Review/New boundary shard.}  This is QSA Step 20 from
\path{docs/quantum_system_associations/PLAN.md}.  It adds no source,
convention entry, script, run bundle, generated data, or analytic
completion.  Its purpose is to mark where the finite association catalogue
touches infinite spectra.

The controlling conventions are \texttt{(v)}--\texttt{(w)} for open-local
observable nets, \texttt{(y)} for closed-point finite-residue sites on
finite-type affine \(\mathbb F_q\)-schemes, \texttt{(z)} for rational-function
Weyl sectors over \(K=\mathbb F_q(t)\), \texttt{(ai)} for the
\(\Spec\mathbb Z\) closed-prime layer, and \texttt{(ak)}--\texttt{(al)} for
regular-function profiles and finite smearings.  The report anchors are
AQM-25, AQM-29, AQM-30, AQM-45, AQM-48, AQM-75, and AQM-79.

\subsection{Boundary Vocabulary}

Every infinite construction below is assigned one of the following statuses.

\begin{description}
\item[Finite-support.] A finite subset \(S\) of the selected closed points is
  quantized by an ordinary finite tensor product of local matrix algebras.
\item[Algebraic quasi-local.] The local algebra is the filtered colimit over
  finite closed supports inside an open set.  No Hilbert-space completion or
  \(C^*\)-completion is part of this status.
\item[Incomplete tensor representation.] A Hilbert representation is chosen by
  fixing a reference vector at every closed point, so basis vectors have
  finite excitation support relative to that reference.
\item[Formal infinite-volume.] A regular-function profile or similar global
  label has infinite closed support.  Its finite restrictions exist, but the
  full profile is not an element of the algebraic quasi-local Weyl algebra.
\item[Generic-sector proposal.] A generic or curve-generic residue field is
  given a separate algebraic Weyl sector only after a character and
  representation convention are stated.  It is not a finite closed-point site.
\item[Completion proposal.] A local-field, adelic, profinite, analytic, or
  \(C^*\)-completed object is named as a possible later layer, not as a
  present report construction.
\end{description}

This vocabulary is only a status partition for Step 21.  It does not assert
agreement or inequivalence between workflows.

\subsection{\texorpdfstring{\(\Spec\mathbb F_q[t]\)}{Spec Fq[t]}}

AQM-29 is the affine-line boundary example.  The scheme has the generic point
\(\eta=(0)\), with residue field \(\mathbb F_q(t)\), and closed points
\[
  x_\pi=(\pi)
\]
indexed by monic irreducible polynomials.  A closed point of degree \(d\)
has residue field \(\mathbb F_q[t]/(\pi)\) and contributes one
\(q^d\)-level finite-field Weyl site.

The finite-support construction is the finite tensor product over a selected
finite set \(S\) of closed points:
\[
  \mathcal A(S)=
  \bigotimes_{x_\pi\in S}
  \operatorname{End}_{\mathbb C}\ell^2(\mathbb F_q[t]/(\pi)).
\]
For an open \(U\), the algebraic quasi-local construction is
\[
  \mathcal A(U)=
  \varinjlim_{S\Subset U\cap |X|_{\rm cl}}\mathcal A(S).
\]
This is the status of the closed-point field net.  It is not a single finite
Hilbert space.

AQM-48 records an incomplete tensor representation in the case
\(q=3\): after choosing \(|0\rangle_\pi\) at every closed point, the
configuration basis is indexed by
\[
  \Gamma_{\rm fin}=\bigoplus_\pi \mathbb F_3[x]/(\pi).
\]
This Hilbert-space representation is extra data.  It is not forced by
\(\Spec\mathbb F_3[x]\) alone.

AQM-29 and convention \texttt{(ak)} also record the formal infinite-volume
boundary.  A nonzero polynomial pair gives a closed-point residue profile
with infinite support.  Its finite restrictions define finite-volume Weyl
operators, but the full regular-function profile is not in the algebraic
quasi-local algebra until a separate infinite-product or completion convention
is fixed.

The generic point is not a finite Weyl site.  Convention \texttt{(z)} gives a
separate rational-function Weyl layer for \(K=\mathbb F_q(t)\) only after an
additive character and representation convention are chosen.  The LCA/local
completion alternatives in \texttt{(z)} remain completion proposals here.

\subsection{\texorpdfstring{\(\Spec\mathbb F_q[x,y]\)}{Spec Fq[x,y]}}

AQM-30 is the affine-plane boundary example.  Its point types are the generic
point, curve-generic points \((f)\), and closed maximal ideals.  Only the
closed maximal ideals are finite-residue Weyl sites under convention
\texttt{(y)}.  If a closed point \(\mathfrak m\) has residue degree
\(d(\mathfrak m)\), then the local carrier is
\[
  \ell^2(\kappa(\mathfrak m)),
  \qquad |\kappa(\mathfrak m)|=q^{d(\mathfrak m)}.
\]

Finite-support and algebraic quasi-local statuses are exactly as in the
affine-line case, with finite subsets of
\[
  U\cap\operatorname{Max}\mathbb F_q[x,y].
\]
The rational-point support
\[
  \{(x-a,y-b):a,b\in\mathbb F_q\}
\]
is a finite-support truncation with \(q^2\) \(q\)-level sites.  It is not the
full \(\Spec\mathbb F_q[x,y]\) theory, because higher-degree closed points are
omitted.

Coordinate profiles such as \((x,0)\) and \((0,y)\) have formal
infinite-volume status in AQM-30.  Their finite restrictions are legitimate
finite-support labels, but the full profiles are not quasi-local Weyl
operators.

The generic point and curve-generic points are generic-sector or completion
questions, not finite Weyl sites.  Their residue fields are rational-function
fields or function fields of curves.  No local-field completion,
Pontryagin-dual phase space, adelic replacement, \(C^*\)-completion, or
generic-point Hilbert-space convention is fixed for them in this shard.

\subsection{\texorpdfstring{\(\Spec\mathbb Z\)}{Spec Z}}

AQM-45 and convention \texttt{(ai)} give the integer boundary example:
\[
  X=\Spec\mathbb Z,\qquad \eta=(0),\qquad x_\ell=(\ell).
\]
The closed prime \(x_\ell\) has residue field \(\mathbb F_\ell\), so the
finite-support closed-prime layer uses
\[
  \mathcal A(S)=
  \bigotimes_{\ell\in S}M_\ell(\mathbb C)
\]
for a finite set \(S\) of rational primes.

For an open \(U\), the algebraic quasi-local closed-prime algebra is
\[
  \mathcal A_{\rm cl}(U)=
  \varinjlim_{S\Subset U\cap\{x_\ell\}}\mathcal A(S).
\]
This is the algebraic infinite tensor product over finite prime supports when
\(U=X\).  It is not a canonical Hilbert-space representation.

AQM-45 also records an incomplete tensor representation after choosing
\(|0\rangle_\ell\in\ell^2(\mathbb F_\ell)\) at every prime.  The basis is
indexed by finite-support residue configurations
\[
  \bigoplus_\ell \mathbb F_\ell.
\]
This representation is useful for local finite-excitation calculations, but
it is not a \(C^*\)-completion and it does not implement every profinite-type
automorphism by a unitary in that sector.

The generic point has residue field \(\mathbb Q\).  AQM-45 records an optional
algebraic \(\mathbb Q\)-Weyl sector with carrier \(\ell^2(\mathbb Q)\), but it
is present in every nonempty open and is a global generic-fibre layer, not a
prime-local qudit.  Adelic, local-completion, and analytic Stone--von Neumann
versions remain completion proposals until separate conventions and sources
are added.

\subsection{Representation and Completion Gaps}

The missing choices for infinite boundary cases are now explicit.

\begin{itemize}
\item Infinite residue fields do not inherit the finite closed-point
  \(\ell^2(\kappa(x))\) qudit convention.  They need a character, dual-label
  space, and representation convention.
\item Generic points and curve-generic points do not become finite Weyl sites
  by default.  If included, they are generic-sector proposals.
\item Analytic or local completions require locally compact groups,
  Pontryagin duals, continuity hypotheses, and a sourced Stone--von
  Neumann--Mackey statement before uniqueness language is allowed.
\item \(C^*\)-completions of the algebraic quasi-local algebras require a
  norm, completion, and state-continuity convention.
\item Adelic and profinite sectors require their own topology, restricted
  product, measure, and representation conventions.  The profinite residue
  action in AQM-45 is not the same as an adelic Hilbert-space model.
\item Formal infinite-volume regular-function profiles require a completion
  or infinite-product field-operator convention before they become operators.
\end{itemize}

\subsection{Step 21 Handoff}

Step 21 may use this shard as a status classifier.  Rows for
\(\Spec\mathbb F_q[t]\), \(\Spec\mathbb F_q[x,y]\), and \(\Spec\mathbb Z\)
should distinguish finite-support truncations, algebraic quasi-local
closed-point algebras, incomplete tensor representations, formal
infinite-volume profiles, generic-sector proposals, and completion proposals.
This shard deliberately does not decide final agreement, inequivalence, or
degeneracy claims.  Those require the explicit maps, invariants, or
obstructions demanded by the Step 21 synthesis.
